% Author-review draft. See EDITORIAL.md for the remaining submission work.
% The 2026 Haskell Symposium format is a planning baseline, not a target date.
\documentclass[acmsmall,screen,nonacm]{acmart}
\usepackage{listings}
\usepackage{booktabs}
\usepackage{tabularx}
\settopmatter{printacmref=false}
\setcopyright{none}
\citestyle{acmauthoryear}
\lstset{language=Haskell,basicstyle=\small\ttfamily,
  columns=fullflexible,keepspaces=true,showstringspaces=false,
  commentstyle=\normalfont\ttfamily,breaklines=true,frame=single,
  xleftmargin=0.4em,xrightmargin=0.4em,aboveskip=0.7em,belowskip=0.7em}
\newcommand{\code}[1]{\texttt{#1}}
\newcommand{\sbv}{\textsc{SBV}}
\newcommand{\sourceversion}{14.8}
\newcommand{\paperlisting}[1]{\par\noindent
  \begin{minipage}{\linewidth}
  \lstinputlisting{build/snippets/#1.hs}
  \end{minipage}\par}
\newcommand{\authorplaceholder}[2]{\par\medskip\noindent
  \fbox{\begin{minipage}{0.94\linewidth}\small
  \textbf{Author material to add: #1.} #2
  \end{minipage}}\par\medskip}

\title{Verification as a Library: Lessons from Sixteen Years of SBV}
\subtitle{Experience Report}
\author{Levent Erkok}
\authorsaddresses{}
\renewcommand{\shortauthors}{Levent Erkok}

\begin{document}

\begin{abstract}
SMT solvers make many verification problems accessible to programmers, but
turning a program into a faithful solver query remains a substantial engineering
task. SBV addresses this task as a Haskell library: programmers construct typed
symbolic computations, state properties, and inspect counterexamples using
ordinary Haskell abstractions. Developed since 2010, the library has grown from
bit-vector verification to incremental solver interaction, symbolic algebraic
datatypes, structured inductive proofs, and C generation. This report examines
the architectural decisions and semantic boundaries exposed by that evolution.
Three lessons recur: a symbolic embedding must make the boundary between host
evaluation and logical computation explicit; incremental interaction makes
declaration lifetime and dependency discovery part of the semantics of the API;
and multiple interpretations of one symbolic representation require independent
validation at their boundaries. We develop these lessons through executable
examples, documented release changes, and the current regression infrastructure.
The resulting account identifies both the reach of verification as a library
and the obligations that remain with the programmer, the embedding, and the
solver. It offers design guidance for Haskell libraries that connect a typed,
functional interface to a stateful reasoning engine.
\end{abstract}

\keywords{Haskell, symbolic programming, SMT, embedded languages, verification, experience report}
\maketitle

\section{Introduction}
\label{sec:introduction}

An appealing way to introduce verification to a Haskell programmer is to start
with a function and ask whether a property holds for every input. The notation
can be almost familiar:
\begin{lstlisting}
prove $ \x -> x + 1 .> (x :: SInteger)
-- Q.E.D.

prove $ \x -> x + 1 .> (x :: SInt8)
-- Counterexample: x = 127
\end{lstlisting}
The first claim concerns unbounded mathematical integers. The second concerns
eight-bit signed arithmetic, where the successor of 127 is -128. The change of
type changes the theorem. Neither a pleasant interface nor a powerful solver
removes the need to choose the right semantics.

\sbv{}~\citep{sbv} provides symbolic counterparts of Haskell types and operations
for expressing such problems. A program over symbolic values constructs a
representation that can be translated into SMT-LIB~\citep{smtlib}, submitted to
an external solver, and interpreted as a proof result or a model. The same
infrastructure also supports constraint solving, optimization, interactive
queries, and compilation of supported computations to C. Haskell supplies the
host language, the type system, and much of the program-construction machinery.

The library began in 2010. Its release history records a first stable public
Hackage release in January 2011, early C generation and changes to sharing,
an incremental execution model in 2017, transformer support in 2019, and major
extensions to proofs and datatypes in 2025--2026~\citep{sbvchanges}. This history
provides an opportunity to examine a question that is broader than any one
feature: \emph{which boundaries must an embedded verification library preserve
as its capabilities grow?}

The answer developed here concerns three boundaries. First, Haskell evaluates
the program that constructs a symbolic problem; symbolic operations describe
the problem itself. Second, a typed expression exists within a particular
symbolic run and, during a query, a particular evolving solver session. Third,
concrete evaluation, SMT interpretation, and generated code must agree wherever
they claim to implement the same operation. Several apparently unrelated
maintenance problems become understandable as failures to preserve one of
these boundaries.

This report contributes an architectural account of the typed embedding and
its shared representation; an analysis of documented changes to sharing,
incremental execution, and datatype registration; and executable case studies
connecting automated checking, inductive proof, and C generation. The
contribution is experience with composing and maintaining these mechanisms.
We do not claim a new SMT decision procedure, a new general technique for
observable sharing, or a verified implementation of the library.

Our evidence is the \sourceversion{} source snapshot, its release history,
selected regression tests, and the small reproducibility companion supplied
with this draft. The scope is deliberately explicit: examples demonstrate
particular capabilities, while historical changes expose particular engineering
pressures. Neither is a measurement of adoption or a comparative performance
study. Section~\ref{sec:evidence} states the limits of this evidence.

\section{A Small Program with Several Interpretations}
\label{sec:example}

Consider addition of unsigned bytes that saturates at 255 rather than wrapping.
The implementation can test for overflow before adding. Its specification can
use unbounded integers, making the intended mathematical result explicit:
\paperlisting{saturating}
The conditional \code{ite} and equality \code{.==} operate on symbolic values.
The specification converts each byte to an unbounded integer before adding.
The expression \code{255 - y} does not
underflow because every \code{SWord8} value is in the range $0$ to $255$.
\code{prove satAddSpec} asks whether the two expressions agree for every pair of
byte inputs. A successful result covers the finite input domain through the
solver's reasoning; it is not a sample of concrete executions.

For literal inputs, operations can evaluate during construction. For example,
\code{satAdd 250 10} yields the literal byte 255. For symbolic inputs, the
operations build a graph. For C generation, the same definition supplies the
body of a generated function:
\paperlisting{codegen}
The generator declares the input interface and the returned value. It does not
invoke a solver at C run time. Moreover, proving \code{satAddSpec} does not
establish correctness of the C compiler. The companion therefore also compiles
the generated code and compares all 65,536 input pairs with an independently
written C oracle using wider unsigned arithmetic. The theorem and the exhaustive
execution check establish different facts.

Even a minimal example has several contracts. The input types specify the
arithmetic domain. The property specifies the desired behavior. The embedding
must encode both correctly. The solver must answer the encoded question
correctly. The code generator must implement the symbolic operations correctly
under its documented C assumptions. Separating these contracts makes it possible
to use one representation productively without attributing every guarantee to
every consumer of that representation.

The familiar arithmetic syntax is selective. Haskell's \code{Num} class allows
overloading of addition and literals. Its \code{Eq} class returns \code{Bool},
so a symbolic equality needs a distinct operation returning \code{SBool}.
Similarly, ordinary Haskell \code{if} inspects a concrete boolean; it cannot
choose between branches on a symbolic condition. The resulting interface mixes
ordinary operators with symbolic combinators. This visible distinction helps
locate the boundary between constructing a problem and asking a logical
question about it.

\section{The Representation Behind the Interface}
\label{sec:architecture}

\subsection{Typed values, concrete values, and graph nodes}

The central public type is \code{SBV a}. For example, \code{SBool} is
\code{SBV Bool}, \code{SInteger} is \code{SBV Integer}, and \code{SWord8} is
\code{SBV Word8}. Internally, the representation has the following shape; this
listing omits fields and constructors not needed for the discussion:
\begin{lstlisting}
newtype SBV a = SBV { unSBV :: SVal }
data SVal = SVal Kind (Either CV (Cached SV))
data SV = SV Kind NodeId
data SBVExpr = SBVApp Op [SV]
newtype SymbolicT m a =
  SymbolicT { runSymbolicT :: ReaderT State m a }
\end{lstlisting}
The parameter of \code{SBV} is phantom in this representation but carries the
public type information. \code{Kind} records the symbolic sort at run time;
\code{CV} represents a concrete value. A non-concrete value contains a deferred
computation that produces a symbolic node within a state. \code{SymVal} connects
Haskell types to symbolic kinds and concrete representations.

This separation avoids making every internal graph operation a heterogeneous
Haskell term indexed by its full type. It also creates implementation
obligations: the phantom interface, run-time kinds, concrete payloads, and
operation signatures must agree. The public API rules out many ill-typed
combinations, but it is not a type-level proof of the implementation. Internal
checks remain necessary, especially for casts and dynamically represented
operations.

\begin{figure}[t]
\centering
\fbox{\begin{minipage}{0.91\linewidth}
\centering
Haskell program over \code{SBV a}\\[3pt]
$\downarrow$\quad concrete evaluation and symbolic construction\\[3pt]
Shared graph, kinds, constraints, and function definitions\\[5pt]
\begin{tabular}{ccc}
SMT-LIB translation & & C lowering\\
$\downarrow$ & & $\downarrow$\\
Solver session & & C source and runtime helpers\\
$\updownarrow$ & & $\downarrow$\\
Queries and models & & Executable program
\end{tabular}\\[5pt]
TP constructs proof obligations using the symbolic and query interfaces.
\end{minipage}}
\caption{The shared representation and its consumers. The arrows describe
implementation flow, not a proved correctness relation between backends.}
\Description{Haskell constructs a shared symbolic graph. One branch translates
the graph to SMT-LIB and connects a solver to queries and models. Another
lowers it to C source and runtime helpers. TP uses the symbolic and query
interfaces to construct proof obligations.}
\label{fig:architecture}
\end{figure}

Concrete evaluation is embedded in the operation constructors. When arguments
are known, an operation may produce a concrete result; otherwise it records an
application of an operation to symbolic nodes. Algebraic simplifications can
avoid graph construction as well. These simplifications must respect the
operand kind. In particular, identities justified for integers cannot simply
be reused for IEEE floating-point values, whose semantics includes NaNs, signed
zero, and rounding.

\subsection{Two forms of sharing}

The implementation preserves sharing at two levels. Deferred symbolic
computations are memoized within a symbolic state using stable names. Separately,
\code{newExpr} consults a map of symbolic applications to reuse existing nodes.
It normalizes certain argument orders and checks that a reused node has the
requested kind. The latter check matters because an expression can participate
in signedness-changing operations.

These mechanisms solve related but different problems. Reusing the result of a
deferred computation avoids repeatedly constructing the same graph fragment.
Reusing an existing application identifies equal constructed expressions even
when their Haskell computations were separately allocated. Neither mechanism
turns the graph into a canonical representation of mathematical equivalence:
two equivalent formulas can still have different nodes.

\authorplaceholder{a design decision}{Describe one consequential choice in
the embedding: the alternatives considered, the reason for the choice, and
what subsequent experience changed. Sharing is a candidate, but the author
should choose the episode. Do not infer personal motivations from the code.}

The state contains mutable tables, counters, and accumulated assignments.
\code{SymbolicT} presents this environment through a reader transformer, with
mutation performed through references. \code{Symbolic} specializes the
transformer to \code{IO}. The transformer interface allows clients to combine
symbolic construction with their own effects; the ordinary interface keeps
simple uses short. The 2019 release introduced this transformer support while
retaining the usual entry points~\citep{sbvchanges}.

\subsection{Values belong to runs}

Haskell type equality does not establish that two symbolic nodes belong to the
same symbolic execution. Nodes therefore carry context information, and graph
construction checks for common context mismatches. Returning a symbolic value
from one solver run and combining it with a value from another can violate this
invariant even when both have type \code{SInteger}.

This is a limitation of what the current public types express. In contrast,
extracting a concrete \code{Integer} from a model and using it to construct a
literal in a new run has a clear interpretation. The distinction suggests a
general rule for embedded libraries: sharing the host type of a value does not
establish the lifetime of the external resource to which that value refers.
The implementation's context checks detect common violations, rather than
providing a static guarantee for every possible misuse.

\section{Evolution Exposes the Boundaries}
\label{sec:evolution}

Table~\ref{tab:history} selects changes relevant to the argument. The dates and
descriptions are taken from the release history; the lessons below are an
analysis of those changes and their current implementations.

\begin{table}[t]
\caption{Selected milestones recorded in the release history~\citep{sbvchanges}.}
\label{tab:history}
\begin{tabularx}{\linewidth}{@{}l l X@{}}
\toprule
Date & Release & Change relevant to this report\\
\midrule
2010 & 0.0.0--0.9.6 & Initial infrastructure and design exploration\\
2011-03 & 0.9.12, 0.9.14 & C generation; sharing reworked using stable names\\
2017-07 & 7.0 & Interleaved symbolic execution and solver interaction\\
2019-01 & 8.0 & Symbolic computations integrated with transformer stacks\\
2025-02 & 11.1 & Reworked calculational and inductive proof interfaces\\
2025-10 & 13.0 & User-defined symbolic algebraic datatypes\\
2026-09 & 14.8 & New C backend; query and datatype registration fixes\\
\bottomrule
\end{tabularx}
\end{table}

\subsection{Host evaluation is part of the embedding design}

The March 2011 release history records replacing an earlier sharing mechanism
with one based on stable names, inspired by \code{Data.Reify}. It also records
problems involving constant applicative forms in the earlier approach. This is
a concrete instance of the observable-sharing problem studied by
\citet{sharing}: a functional description can share a computation, while a
naive reification expands it into a much larger tree.

For \sbv{}, the practical requirement is to recover reusable structure without
making the ordinary arithmetic interface monadic at every operation. The
current design defers node creation until a symbolic state is available and
memoizes within that state. The historical record supports the need to revise
this mechanism; it does not by itself quantify the performance improvement.

Recursion reveals the same host/object distinction from a different direction.
A Haskell function that recursively calls itself under \code{ite} may continue
constructing symbolic branches indefinitely. A finite input type does not
guarantee a finite construction: unless the stopping condition becomes concrete
during host evaluation, the construction can keep unfolding. The Fibonacci
code-generation example in the repository illustrates using an explicit
construction bound to obtain a finite graph.

Named symbolic functions provide another route. \code{smtFunction} represents
a function definition and its calls at the SMT level, allowing recursive
definitions to remain recursive rather than being fully unfolded by Haskell.
That changes where the recursion is represented; it also introduces obligations
about definitions and termination. The library can help discharge those
obligations, but moving a recursive call into the solver is not itself a proof
that the function is well founded.

\subsection{Incremental solving changes the lifetime of declarations}
\label{sec:incremental}

The 2017 release documents a change from running a symbolic program before
invoking the solver to interleaving execution with solver interaction. Its
stated motivation was to let clients send and receive solver commands instead
of being limited to predefined recipes. The earlier tactic interface was
removed as part of that change~\citep{sbvchanges}.

The query interface makes it possible to inspect a model and decide what to ask
next. The following example enumerates the integer roots of a polynomial:
\paperlisting{roots}
After a satisfiable result, the program obtains a concrete root and adds a
constraint excluding it. The next check runs in the same session. Enumeration
ends when the remaining constraints are unsatisfiable. An inconclusive solver
response is handled separately; it is not treated as exhaustion. Root order is
not part of the contract.

Such a small loop has consequences for the implementation. Literals,
expressions, types, and functions can be encountered after the initial problem
has been sent. The library must record newly created objects and synchronize
them with the solver before a command refers to them. Its symbolic state
therefore maintains an incremental component as well as the full accumulated
program. Query operations use that component to emit new declarations,
definitions, assignments, and constraints.

Release 14.8 fixes first use of \code{smtFunction} definitions during queries
without prior explicit registration. It also fixes first use of rational kinds
in query mode, including occurrences nested inside other types. These changes
demonstrate that incremental support is not complete merely because
\code{checkSat} and \code{getValue} are available. Every operation that can
introduce a dependency must participate in synchronization.

The resulting lesson is about API compositionality. If an expression can be
constructed during setup and during a query, its dependencies should normally
follow it in either context. Requiring an unrelated setup-time use to make a
later expression work exposes an implementation accident. Some solver options
genuinely must be chosen before a session starts; those restrictions should be
distinguished from missing dependency registration.

\authorplaceholder{a historical surprise}{Add a specific user request, solver
behavior, or failed design that changed the direction of SBV. Include an
approximate date and, where available, a release or issue reference. Explain
the resulting lesson rather than only recounting the event.}

\subsection{Datatypes make dependencies structural}

\sbv{} uses Template Haskell to make user-defined datatypes symbolic. A splice
such as \code{mkSymbolic [''Box]} generates the symbolic infrastructure for a
Haskell datatype. \code{sCase} provides symbolic pattern matching, including
nested patterns and guards. The mechanism works with the embedding's symbolic
representation; it does not change the meaning of arbitrary Haskell
\code{case} expressions throughout a program.

A datatype dependency need not occur immediately in a constructor field. It
can appear inside a tuple, a list, a set, an array, or a type argument. Mutually
recursive definitions add cycles. Consequently, registering only the visible
root constructor is insufficient. The library must discover the structural
dependencies needed to declare and use the symbolic kind.

The 14.8 release records fixes for precisely these cases. The associated
registration tests vary both the position of a nested datatype and the phase
in which a symbolic value is first created: setup, a named query variable, or
an unnamed query variable. This forms a useful regression matrix. A test of
one nested datatype in a conventional setup block cannot establish that first
use inside a query will work.

An abbreviated example in the companion creates two \code{Box Word8} values
for the first time inside a query and requires them to differ. The result
should be satisfiable. The full repository tests extend that idea to container
nesting, synonyms, and recursive dependencies. The general lesson is to test a
feature at the places where it enters the system, in addition to testing the
feature's operations once it is already registered.

\section{From Automated Checks to Structured Proofs}
\label{sec:proofs}

Push-button checking is effective when the solver can discover the necessary
argument. Recursive functions often need structure supplied by the user:
an induction principle, a stronger invariant, a case split, or intermediate
equalities. \sbv{}'s TP layer lets that structure be written as a Haskell
program that generates and checks proof obligations.

An accumulator-based list reversal illustrates the interaction. Its useful
invariant quantifies over both the input list and the accumulator:
\paperlisting{reverse}
The function is specialized to symbolic lists of integers here to keep the
listing small. \code{Forall} names quantified parameters at the type level.
\code{induct} supplies the list induction structure, including a base case and
an induction hypothesis. The user gives the calculation for the step case.
The operators \code{|-}, \code{=:}, and \code{??} introduce assumptions,
successive expressions, and proof hints respectively.

The crucial human contribution is the generalized statement involving
\code{acc}. A claim only about reversing into an empty accumulator is too weak
to apply directly to the recursive call. Once the stronger statement is
chosen, the calculation exposes the useful intermediate expressions and gives
the solver smaller obligations. The notation expresses a proof strategy;
Haskell's type checker alone does not certify the displayed equalities.

In ordinary checked execution, TP's proof-step machinery negates the requested
proposition under the current assumptions and asks for satisfiability.
Unsatisfiability discharges the step. A satisfiable or inconclusive result does
not. The layer records proof dependencies and supports reuse of lemmas.
Recursive definitions additionally require attention to their termination
measures. The library offers measure-based checks and records uses of functions
whose termination checks were explicitly bypassed.

The constant-folding development is a larger example already shipped with the
library. Its expression datatype includes variables, constants, arithmetic,
and local bindings. The transformation recursively folds children and then
simplifies their parent. Replacing a binding by a constant requires
substitution that respects shadowing. Its main claim is:
\[
  \forall e,\rho.\quad
  \operatorname{interpInEnv}(\rho,\operatorname{cfold}(e))
  = \operatorname{interpInEnv}(\rho,e).
\]
The proof uses a size measure and auxiliary lemmas about substitution,
environments, and arithmetic. This example makes an important limitation
visible: the convenience of an embedded prover does not eliminate the work of
discovering the supporting theory. The paper companion invokes the existing
development; it does not claim a new constant-folding algorithm or a new proof
of that algorithm.

These are claims about the represented semantics. \code{SList a} denotes a
finite sequence of values of type \code{a}; it is not the domain of arbitrary
lazy Haskell lists with partial and infinite values. A reversal theorem over
symbolic lists therefore does not establish the corresponding statement for
all Haskell list computations. Likewise, an interpreter represented with
symbolic datatypes has the semantics of its symbolic definitions. Relating it
to a separately implemented interpreter is a further verification task.

\section{Code Generation Tests the Common Representation}
\label{sec:codegen}

C generation has been present since 2011, but the 14.8 backend substantially
extends its scope. It supports, among other values, arbitrary-width
bit-vectors, exact integers, rational-valued reals, floating-point values,
structured values, and recursive function definitions. Exact unbounded
integers and rational-valued reals use GMP; some floating-point operations
require LibBF. The legacy backend remains separately available.

Sharing the expression graph avoids repeating the front-end representation,
but the target has different operational requirements. A solver reasons about
mathematical terms. C needs a choice of evaluation order, storage, ownership,
and failure behavior. The backend therefore lowers an expression to a result
expression together with declarations, setup statements, and required runtime
facilities. Setup statements must execute when their values are needed. Moving
them across conditional branches indiscriminately can change behavior.

Numeric operations expose another boundary. The release notes record fixes to
pseudo-Boolean comparisons whose totals exceed the range of a naive native
accumulator, and to flooring of native-mapped reals when converting to
integers. Both operations can be easy to express symbolically while being
subtle to implement in C. The regression suite uses independent mathematical
expectations and executes generated code, including checks at different
optimization levels and with sanitizers in selected tests.

The generator has explicit limits. Quantifiers and other solver-only
constructs are rejected. Exact real support covers rational values rather
than arbitrary algebraic reals. Direct array equality is supported only for
enumerable finite key domains within a configured budget. Floating-point
generation has documented rounding and compiler-option requirements. These
boundaries are part of the meaning of using the backend, not implementation
details that can be hidden behind a common API.

The saturating-adder case in Section~\ref{sec:example} is intentionally small
enough to validate exhaustively. It shows one way to test the connection between
the specification, symbolic implementation, and generated code. It supplies
no blanket guarantee for the other operations or runtime representations.

\section{Evidence and Reproducibility}
\label{sec:evidence}

The source snapshot used here is commit
\code{f4f7a7a08445e63c2fb74c6f2084dfea5ad8b7d2}, identifying release 14.8.
The companion's \code{README.md} records build commands and the local tool
versions; \code{results/validation.txt} records the executed checks. The paper's
main executable listings are extracted from the checked Haskell source rather
than maintained as separate copies.

\subsection{Small, explicit validation tasks}

Table~\ref{tab:checks} defines the companion's tasks and the evidence each
supplies. All nine tasks completed successfully on 26 September 2026 using
GHC 9.14.1, Z3 5.1.0, and CVC5 1.4.0 on an ARM64 macOS host. The generated C
was compiled with Apple Clang 21.0.0, optimization enabled, strict warnings,
and undefined-behavior sanitization. The checks cover distinct paths through
the library. Their
small scale is deliberate: they make the claims in the text reproducible
without turning an experience report into a broad solver benchmark.

\begin{table}[t]
\caption{Companion validation tasks. These are semantic checks, not timing
benchmarks; the execution record accompanies the draft.}
\label{tab:checks}
\begin{tabularx}{\linewidth}{@{}l X@{}}
\toprule
Task & Required outcome\\
\midrule
Integer successor & Theorem holds over unbounded integers\\
Signed-byte successor & Counterexample is the value 127\\
Floating-point reflexivity & A NaN is returned as a counterexample\\
Saturating addition & Implementation satisfies the integer specification\\
Incremental roots & Exactly $-3$ and $3$, followed by unsatisfiability\\
Late datatype declaration & Distinct query-created byte boxes are satisfiable\\
Accumulator reversal & TP completes the generalized inductive proof\\
Constant folding & The existing CVC5 proof development completes\\
Generated C addition & All 65,536 input pairs agree with a wider-arithmetic oracle\\
\bottomrule
\end{tabularx}
\end{table}

The companion checks model values rather than assuming a particular printed
model order. It treats timeouts as failures to complete a task, not as evidence
that a proposition is false. It also checks the generated adder against an
oracle whose arithmetic domain is wider than the inputs. These choices keep
the evidence close to the semantic claim being made.

\subsection{Regression infrastructure as experience evidence}

The repository's testing infrastructure combines HUnit tests, QuickCheck
properties, golden files, compile-time diagnostics, executable documentation,
and generated-C integration tests. Each addresses a different failure mode.
The shared \code{qc1} and \code{qc2} helpers, for example, can exercise an
operation both on concrete symbolic literals and through a solver. This is a
direct response to the need for concrete evaluation and SMT translation to
agree.

Golden files make changes to models, proof transcripts, and generated text
reviewable. They are valuable records of behavior, but their evidential force
depends on what was recorded and reviewed. A stable transcript is not an
independent proof that the translation was correct. Similarly, compile-failure
tests for \code{sCase} and \code{pCase} exercise the diagnostics and rejection
behavior of the metaprogramming interface; they address concerns that a
successful solver example cannot test.

The datatype-registration matrix and executed C boundary tests are especially
useful examples of tests derived from architectural obligations. They vary the
phase, nesting, representation, or compilation mode that previously exposed a
problem. That is more informative than demonstrating a newly added operation
only in an already well-prepared context.

\subsection{Limits of the account}

\authorplaceholder{an application outside the example suite}{Describe a real
use of SBV: the problem, why this approach was chosen, what worked, what did
not, and evidence that can be cited or reproduced. Identify any limitations
on naming the application. This material will extend the present account
beyond examples and maintenance records.}

The examples and regressions were selected to explain the three boundaries
identified in Section~\ref{sec:introduction}. They are not a random sample of
all uses of \sbv{}, and they do not estimate failure rates. The release notes
provide contemporaneous records of changes, but not a controlled study of
their causes or effects. We make no measured claim that the API reduces
development effort, that every supported solver behaves identically, or that
the current implementation is free of semantic bugs.

The source contains many larger examples, but the existence of an example is
not evidence that arbitrary problems in the same domain will be tractable.
Quantifiers, nonlinear arithmetic, and recursive definitions can make solver
behavior difficult to predict. Different solver releases can require different
proof guidance. The successful small checks establish only their stated
outcomes in the recorded environment.

\section{Related Work}
\label{sec:related}

\paragraph{Properties as Haskell programs.}
QuickCheck~\citep{quickcheck} established an influential style in which
programmers express properties in Haskell and use generated inputs to seek
counterexamples. \sbv{} shares the appeal of a library-level interface but
constructs logical formulas for supported symbolic operations. Testing and
SMT reasoning offer different forms of evidence and can be used together, as
they are within \sbv{}'s own regression suite. The syntax of a property alone
does not determine its evidential strength.

\paragraph{Reification and sharing.}
\citet{sharing} develops type-safe observable sharing for embedded languages.
\sbv{}'s release history explicitly connects its stable-name implementation
to \code{Data.Reify}. Our account concerns how deferred construction,
expression reuse, and per-run state coexist in this particular library. It
does not present observable sharing as a new technique.

\paragraph{High-level symbolic computation.}
Smten~\citep{smten} translates high-level symbolic computations into SMT
queries, addressing the burden of manually constructing effective formulas.
Rosette's symbolic virtual machine~\citep{rosette} supports solver-aided host
languages and uses symbolic evaluation to make rich program descriptions
available to solvers. These systems share the goal of raising the level at
which a user describes a reasoning problem. The relevant distinction here is
the boundary exposed by a Haskell library whose symbolic values and operations
are explicit in the host program. This report does not attempt a feature or
performance ranking of these systems.

Grisette~\citep{grisette} is a particularly close comparison: it provides
symbolic compilation as a typed functional library in Haskell, with a purely
functional representation and an ordered-guards approach to state merging.
Its work demonstrates that a library embedding is itself an established
research direction. \sbv{}'s contribution here is a longitudinal account of a
different engineering composition: deferred graph construction with per-run
state, incremental SMT sessions, structured proof orchestration, and a C
consumer of the symbolic representation. A direct performance comparison would
require carefully matched workloads and is outside this report.

\paragraph{SMT-assisted deductive proofs.}
Refinement reflection~\citep{reflection} connects user-defined functions to
SMT reasoning within Liquid Haskell's refinement-type verification framework.
TP instead presents explicit proof programs whose obligations are discharged
through \sbv{}'s symbolic and query interfaces. Both approaches let users
supply structure that automated reasoning alone may not discover. Their
relationship calls for a comparison of specification styles, trusted
components, and proof guidance, rather than an assumption that all
SMT-assisted proofs provide the same interface or guarantees.

\paragraph{Solvers and their interface.}
SMT-LIB~\citep{smtlib} provides the language through which \sbv{} communicates
with its backends. Z3~\citep{z3} is the default solver, with other solvers
selected through configurations and capability descriptions. \sbv{} relies on
these engines for the logical reasoning. Its responsibilities include choosing
faithful encodings, maintaining the protocol, and interpreting responses
without converting an inconclusive outcome into a proof.

\section{Discussion and Lessons}
\label{sec:discussion}

Verification as a library provides a useful division of labor. Haskell handles
program construction, abstraction, and much of the static discipline. \sbv{}
records symbolic meaning and manages the connection to external reasoning.
Solvers discharge the resulting obligations. The experience examined here
suggests three lessons for similar libraries.

First, expose semantic differences at the point where users make choices.
Separate symbolic equality and branching, distinguish a symbolic sequence from
a Haskell container of symbolic elements, and state the domain of a theorem.
Familiar notation is useful when its interpretation is clear.

Second, treat incremental execution as a property of the whole interface.
Registration, recursive definitions, constants, and nested types must work
when first encountered during a live session. Tests should vary the phase of
first use as systematically as they vary operand types.

Third, use a shared representation to organize independent checks of its
interpretations. Concrete evaluation, SMT translation, and generated C can
cross-check selected operations, but common implementation structure can also
propagate a mistake. Independent mathematical specifications and execution
oracles are therefore valuable at semantic boundaries.

The trusted base remains substantial: the Haskell implementation and its
runtime, the symbolic encoding, proof orchestration and any supplied axioms,
solver behavior, and, where relevant, the C compiler and generated runtime.
TP proof objects record propositions and dependencies; the ordinary workflow
does not establish every result by replaying a certificate in an independent
small kernel. This account should be read with that trust model in mind.

\section{Conclusion}

Sixteen years of \sbv{} development show how a Haskell library can grow from
bit-vector properties to interactive reasoning and structured proofs while
retaining a compact programming interface. The documented changes also show
where the cost of that growth appears: host evaluation and sharing, the
lifetime of solver declarations, and agreement between interpretations.
Making those boundaries explicit yields both a clearer user model and more
targeted tests. Those lessons are applicable to embedded libraries that must
preserve a typed functional interface while coordinating stateful external
tools.

\bibliographystyle{ACM-Reference-Format}
\bibliography{references}
\end{document}
